\section{Discussion}

% Shared section

The evaluation shows that earlier governance timing can improve prevention, but
it does not establish that PBCP is immediately ready for production enforcement.
The discussion therefore focuses on the boundary between what the benchmark
supports today and what would still need validation in operational settings.

\subsection{Limitations}

Although the benchmark captures diverse workload behaviors, it remains a
controlled synthetic environment rather than production telemetry. Real
enterprise workloads may exhibit additional operational variance, organizational
approval constraints, and intervention costs not represented here. The current
results also depend on assumed policy and intervention parameters, the policy
learning loop is still evaluated at limited scale, and the system has not yet
been validated inside a live enterprise enforcement stack. Similarly, the
Exp~3 mismatch-subgroup result is mixed even though overall detector quality
improves, so the current anomaly-detection claim should be read as strong
benchmark evidence rather than as final production generalization.

\subsection{Future Work}

The most direct next steps are calibration against real organizational telemetry,
integration with execution platforms such as Databricks or Kubernetes, stronger
online policy learning, and tighter coupling with production approval and
enforcement stacks. These extensions would test whether the same
workload-alignment signals remain useful once governance decisions interact with
real operators, heterogeneous environments, and organization-specific risk
tolerances.

% Authors: Keerthi Rapolu & Sreeja Katta

\subsection{Anomaly Root Cause Analysis and Prevention Feedback}

The RCA module converts recurring divergence patterns into policy suggestions,
closing the governance loop between runtime monitoring and future pre-billing
intervention. After each generation, \texttt{RootCauseAnalyzer} queries
historical incidents grouped by workload type and incident category ---
over-provisioned requests, idle cluster patterns, runaway execution, and
undeclared token budget --- and synthesizes a \texttt{Policy} object when the
same pattern recurs at least three times within a lookback window. Confidence
scales with pattern frequency up to a ceiling of 0.95, and only policies above
the 0.60 threshold enter the registry.

This mechanism addresses a gap that static pre-billing policy cannot fill.
Preconfigured rules handle known divergence patterns, but novel or low-frequency
patterns are not covered at system initialization. The feedback loop fills that
gap incrementally: workloads that escape pre-billing intervention, behave
anomalously at runtime, and are flagged by IFS contribute to learned policies
that catch the same pattern in subsequent submissions. The Exp~6 convergence
result is partly attributable to this accumulation --- Full PBCP substantially
outperforms the No Phase~3 ablation because learned policies progressively
reduce the baseline rate of uncaught divergence.

Two limitations apply to the current implementation. First, the feedback loop
operates on synthetic incident data, so confidence and frequency estimates may
not reflect the noisier incident streams in real enterprise environments. Second,
the pattern-counting logic does not weight incidents by recency or cost severity,
meaning a cluster of low-cost incidents and a smaller cluster of high-cost
incidents are treated identically. Both limitations are addressable in a
production deployment and represent near-term extensions rather than design
defects.

